\documentclass[10pt,a4paper]{article}
\usepackage[margin=1.75cm]{geometry}
\usepackage{amsmath,amssymb,booktabs,graphicx,tabularx,array,xcolor}
\usepackage[hidelinks]{hyperref}
\graphicspath{{exploration/results/tabular_mdp/two_step_trap/}{../results/tabular_mdp/two_step_trap/}}

\newcommand{\pool}{\mathrm{pool}}
\newcommand{\good}{p_{\mathrm{good}}}
\newcommand{\R}{\mathbb{R}}
\newcommand{\E}{\mathbb{E}}
\newcommand{\grad}{\nabla_{\phi}}
\newcommand{\coloneqq}{\mathrel{\mathop:}=}

\title{Step 3: Exact Two-State Tabular MDP Geometry\\
\large A self-contained reconstruction of the motivation, mathematics,
implementation, experiments, and conclusions}
\author{Vanilla Policy Gradient Project}
\date{August 2026}

\begin{document}
\maketitle

\begin{center}
\fbox{\parbox{0.94\textwidth}{
\textbf{What this document is for.}
This is a reconstruction guide, not only a result summary.  Starting from the
multi-armed-bandit log barrier of Cesani, Papini, and Restelli, it explains why
we introduced a two-state MDP, defines every mathematical object used in the
code, derives the six regularizers, states what each experiment isolates, and
then interprets the generated figures without claiming more than the exact
tabular evidence supports.
}}
\end{center}

\tableofcontents

\section{Scientific path: why Step 3 exists}

\subsection{The problem identified in the bandit paper}

Cesani et al. study a stochastic $K$-armed bandit controlled by a categorical
softmax policy.  Their starting observation is that an ordinary stochastic
policy-gradient update can drive action probabilities close to zero.  Once an
action is extremely unlikely, it is rarely sampled and its score-weighted
reward signal becomes too weak to repair an early mistake.  They describe this
as an exploration failure near the boundary of the probability simplex
\cite[Secs.~1--3]{cesani2026}.  This concern is related to the local nature of
policy-gradient exploration more generally \cite{agarwal2020}.

Their proposed bandit objective is
\begin{equation}
\Phi_\eta(\theta)
=J(\theta)+\frac1\eta\sum_{a=1}^{K}\log\pi_\theta(a),
\label{eq:cesani-objective}
\end{equation}
where $1/\eta$ controls the barrier strength.  For a reduced categorical
softmax policy they prove
\begin{equation}
\log\det F_{\mathrm{red}}(\theta)
=\sum_a\log\pi_\theta(a).
\label{eq:cesani-identity}
\end{equation}
Thus the probability-space log barrier has an exact reduced-Fisher
log-determinant interpretation in the bandit setting
\cite[Sec.~4 and App.~B.3.2]{cesani2026}.  Its gradient is deterministic,
\begin{equation}
\nabla_\theta\sum_a\log\pi_\theta(a)=\mathbf 1-K\pi_\theta,
\label{eq:cesani-gradient}
\end{equation}
so it restores small action probabilities without adding reward noise
\cite[Sec.~5 and App.~B.4]{cesani2026}.

This is geometrically related to natural policy gradient (NPG), but it is not
the same operation.  NPG preconditions the reward gradient by $F^{-1}$;
the log barrier changes the objective so optimization is repelled from the
simplex boundary.  Kakade introduced the policy Fisher as a metric for a
coordinate-aware policy update \cite{Kakade2001}, following the natural-gradient
framework of Amari \cite{Amari1998}.

\subsection{What Steps 1 and 2 established}

The project deliberately used a ladder of increasingly complicated settings:
\begin{enumerate}
\item \textbf{Step 1: exact categorical identity.}  With one categorical
policy, every action was enumerated and
\eqref{eq:cesani-identity}--\eqref{eq:cesani-gradient} were checked in CPU
float64.  There were no rewards, trajectories, or state distribution.
\item \textbf{Step 2: stochastic bandit training.}  Reward noise and sequential
learning were introduced while the categorical Fisher identity remained exact.
This reproduced the barrier's support-preservation mechanism and its
exploration/bias trade-off, corresponding to the experimental questions in
Sections 7 and E of Cesani et al. \cite{cesani2026}.
\item \textbf{Step 3: this experiment.}  We introduce multiple states while
keeping the entire MDP exactly enumerable.  This is the smallest setting in
which the policy affects both an action distribution \emph{and} which state is
represented in a transition-pooled dataset.
\end{enumerate}

\subsection{The question that cannot be answered by a bandit}

A bandit has no state-visitation distribution.  In an MDP, a Fisher constructed
from rollout transitions depends on two policy-controlled objects:
\begin{equation}
\underbrace{\text{which states appear}}_{\text{visitation}}
\quad\text{and}\quad
\underbrace{\pi_\phi(\cdot\mid s)}_{\text{conditional action geometry}}.
\end{equation}
Cesani et al. explicitly state that their exact bandit equivalence does not
directly extend to general RL \cite[Sec.~8]{cesani2026}.  Step 3 therefore asks:
\begin{quote}
When state weights depend on the policy, what is the difference between
averaging statewise categorical barrier gradients, differentiating their
policy-weighted scalar average, protecting states uniformly, protecting only
visitation, and differentiating the pooled Fisher log determinant?
\end{quote}

\section{The exact two-step trap MDP}

\subsection{MDP tuple and transition table}

We use the finite MDP
\begin{equation}
\mathcal M=(\mathcal S,\mathcal A,P,R,\rho_0,H,\gamma)
\end{equation}
with nonterminal states $\mathcal S=\{s_0,s_1\}$, actions
$\mathcal A=\{a_0,a_1,a_2\}$, start distribution
$\rho_0(s_0)=1$, horizon $H=2$, and discount $\gamma=1$.

\begin{center}\small
\begin{tabular}{llll}
\toprule
State & Action & Reward & Next state\\ \midrule
$s_0$ & $a_0$ & $0.5$ & terminal (safe action)\\
$s_0$ & $a_1$ & $0$ & $s_1$ (explore)\\
$s_0$ & $a_2$ & $0$ & terminal (abandon)\\
$s_1$ & $a_0$ & $1$ & terminal (good action)\\
$s_1$ & $a_1$ & $0.2$ & terminal (mediocre action)\\
$s_1$ & $a_2$ & $0$ & terminal (bad action)\\ \bottomrule
\end{tabular}
\end{center}

This MDP is intentionally not a benchmark.  It is a diagnostic construction.
The safe action supplies a local reward of $0.5$ without visiting $s_1$.
The globally better outcome requires both selecting $a_1$ at $s_0$ and then
learning $a_0$ at $s_1$.  A policy that rarely reaches $s_1$ receives little
direct learning signal about its actions, creating a minimal visitation trap.

\subsection{All possible trajectories}

Let $\pi_s(a)=\pi_\phi(a\mid s)$.  There are five complete trajectories:
\begin{center}\small
\begin{tabular}{lll}
\toprule
Actions & Probability & Return\\ \midrule
$(a_0)$ at $s_0$ & $\pi_0(a_0)$ & $0.5$\\
$(a_2)$ at $s_0$ & $\pi_0(a_2)$ & $0$\\
$(a_1,a_0)$ & $\pi_0(a_1)\pi_1(a_0)$ & $1$\\
$(a_1,a_1)$ & $\pi_0(a_1)\pi_1(a_1)$ & $0.2$\\
$(a_1,a_2)$ & $\pi_0(a_1)\pi_1(a_2)$ & $0$\\ \bottomrule
\end{tabular}
\end{center}
Their probabilities sum exactly to one.  Writing
\begin{equation}
q\coloneqq\pi_0(a_1),
\qquad
V_1\coloneqq\pi_1(a_0)+0.2\pi_1(a_1),
\end{equation}
the expected return is
\begin{equation}
\boxed{J(\phi)=0.5\pi_0(a_0)+qV_1.}
\label{eq:return}
\end{equation}
The code independently computes this closed form and the explicit sum over all
five trajectories.

\section{Policy coordinates, scores, and exact reward gradient}

\subsection{Why reduced logits are used}

At each state, action $a_2$ is the reference action.  Four free parameters are
collected in
\begin{equation}
\phi=(\phi_{00},\phi_{01},\phi_{10},\phi_{11})\in\R^4.
\end{equation}
The statewise logits are
\begin{equation}
z_0=(\phi_{00},\phi_{01},0),
\qquad
z_1=(\phi_{10},\phi_{11},0),
\end{equation}
and $\pi_s=\operatorname{softmax}(z_s)$.  Fixing one reference logit removes the
redundant common-logit direction.  The full three-logit Fisher is singular
because adding the same constant to all logits leaves the policy unchanged;
the two-coordinate Fisher is positive definite for every finite $\phi$.
This is the same reduced chart used to establish the determinant identity in
Cesani et al. \cite[Sec.~4]{cesani2026}.

\subsection{Score and categorical Fisher at one state}

For action $a$, the score in reduced coordinates is
\begin{equation}
g_s(a)=\nabla_{\phi_s}\log\pi_s(a)
=e_a^{(2)}-\pi_s(0{:}2),
\label{eq:score}
\end{equation}
where $e_0^{(2)}=(1,0)$, $e_1^{(2)}=(0,1)$, and
$e_2^{(2)}=(0,0)$.  Exact enumeration verifies
\begin{equation}
\E_{a\sim\pi_s}[g_s(a)]=0
\end{equation}
and
\begin{equation}
F_s
=\E_{a\sim\pi_s}[g_s(a)g_s(a)^\top]
=\operatorname{Diag}(\pi_s(0{:}2))
-\pi_s(0{:}2)\pi_s(0{:}2)^\top.
\label{eq:state-fisher}
\end{equation}
This is the standard score-outer-product definition of Fisher information.
Natural-gradient policy methods use such a Fisher as the local policy metric
\cite{Kakade2001}.

\subsection{Exact reward gradient and its REINFORCE interpretation}

For $Q_0=(0.5,V_1,0)$ and $r_1=(1,0.2,0)$, the exact gradient is
\begin{align}
\nabla_{\phi_0}J
&=\pi_0(0{:}2)\odot[Q_0(0{:}2)-J\mathbf 1],\\
\nabla_{\phi_1}J
&=q\,\pi_1(0{:}2)\odot[r_1(0{:}2)-V_1\mathbf 1].
\label{eq:reward-gradient}
\end{align}
The factor $q$ in the second line is crucial: reward provides less direct
learning signal for $s_1$ when that state is rarely reached.  The implementation
also enumerates trajectories and verifies
\begin{equation}
\nabla J
=\sum_\tau p_\phi(\tau)G(\tau)\nabla_\phi\log p_\phi(\tau).
\label{eq:reinforce}
\end{equation}
Equation~\eqref{eq:reinforce} is the exact expectation underlying REINFORCE
\cite{Williams1992} and policy-gradient estimators
\cite{SuttonMcAllester2000}.  Here it has no Monte Carlo variance because all
trajectories are summed.

\section{The transition-pooled state distribution}

\subsection{Definition and derivation}

Define
\begin{equation}
\boxed{\mu_0=\frac1{1+q},\qquad \mu_1=\frac q{1+q}.}
\label{eq:mu}
\end{equation}
We call this $\mu_{\pool}$ and define it precisely as:
\begin{quote}
The population state distribution obtained by uniformly selecting one valid
action transition from an infinitely large transition-pooled rollout dataset.
\end{quote}

Every episode contains exactly one valid $s_0$ transition.  An episode contains
an $s_1$ transition with probability $q$.  Therefore the population expected
transition counts per episode are $(1,q)$; normalizing these counts gives
Equation~\eqref{eq:mu}.  Shorter terminated trajectories contribute fewer rows
to the globally pooled transition dataset.  This matches the averaging
convention used when all valid score rows are concatenated before forming an
outer-product average.

\subsection{Why it is not canonical finite-horizon occupancy}

The standard fixed-horizon time-average occupancy would be
\begin{equation}
d_H(s)=\frac1H\sum_{t=0}^{H-1}\Pr(S_t=s).
\end{equation}
It uses a fixed denominator $H$, whereas $\mu_{\pool}$ normalizes by the number
of \emph{valid} transitions.  These are different estimands.

For a finite batch of $N$ episodes, let $X_i$ indicate that episode $i$ reaches
$s_1$.  The empirical pooled weight is
\begin{equation}
\widehat\mu_1
=\frac{\sum_{i=1}^{N}X_i}{N+\sum_{i=1}^{N}X_i}.
\label{eq:ratio-estimator}
\end{equation}
It has a random denominator.  In general,
$\E[\widehat\mu_1]\neq q/(1+q)$.  Step 3 avoids this finite-batch ratio-estimator
issue by using the population quantity exactly; the next stochastic stage will
measure its bias, variance, and probability of observing zero $s_1$ transitions.

\section{From statewise Fishers to the pooled Fisher}

\subsection{Block structure}

The parameters for $s_0$ and $s_1$ are disjoint.  Under $\mu_{\pool}$ the
transition-pooled Fisher is therefore
\begin{equation}
\boxed{F_{\pool}
=\operatorname{blockdiag}(\mu_0F_0,\mu_1F_1).}
\label{eq:pooled-fisher}
\end{equation}
This exact block structure would generally disappear with shared neural
parameters.  It is a deliberate advantage of the tabular diagnostic.

For a three-action reduced categorical softmax,
\begin{equation}
\det F_s=\prod_{a=0}^{2}\pi_s(a),
\qquad
B_s\coloneqq\log\det F_s=\sum_{a=0}^{2}\log\pi_s(a).
\end{equation}
Because multiplying a $2\times2$ block by $\mu_s$ multiplies its determinant by
$\mu_s^2$,
\begin{equation}
\boxed{
\log\det F_{\pool}
=B_0+B_1+2\log\mu_0+2\log\mu_1.}
\label{eq:pooled-logdet}
\end{equation}

\subsection{Normalized full barrier}

We define
\begin{align}
B_{\mathrm{full}}
&\coloneqq\frac12\log\det F_{\pool}\\
&=\underbrace{\frac12(B_0+B_1)}_{B_{\mathrm{uniform}}}
+\underbrace{\log\mu_0+\log\mu_1}_{B_{\mathrm{visit}}}.
\label{eq:full-decomposition}
\end{align}
The factor $1/2$ is \emph{not canonical}.  It is an action-matched
normalization chosen so that the action component is exactly
$B_{\mathrm{uniform}}$.  Consequently, the difference between the normalized
full barrier and the uniform-state action baseline is exactly the visitation
term.  This makes the experiment interpretable; it is not a claim that all
Fisher log determinants should be divided by two.

\subsection{What the eigenvalues and determinants do and do not say}

$\lambda_{\min}(F_s)$ measures the weakest local categorical direction at state
$s$.  The pooled minimum combines conditional degeneration with the multiplier
$\mu_s$.  A small $\lambda_{\min}(F_{\pool})$ can therefore mean:
\begin{enumerate}
\item the conditional action policy at $s_1$ is collapsing;
\item $s_1$ is rarely represented because $\mu_1$ is small;
\item both.
\end{enumerate}
$\log\det F$ is the sum of log eigenvalues and measures a local volume, not a
condition number.  A larger determinant does not guarantee a larger smallest
eigenvalue or better optimization.  Moreover, the sampled outer-product matrix
often called an empirical Fisher must not automatically be treated as the true
model Fisher or as a Hessian; Kunstner et al. explain this distinction in detail
\cite{Kunstner2019}.

\section{The six exact-gradient objectives}

All methods use explicit Euler ascent
\begin{equation}
\phi_{k+1}=\phi_k+\alpha[\nabla J(\phi_k)+\beta g_B(\phi_k)].
\label{eq:update}
\end{equation}
The reward gradient is identical across methods.  Only $g_B$ changes.

\subsection{A. Reward only}
\begin{equation}
g_B=0.
\end{equation}
This is the control.  It shows the basin selected by the exact reward objective
without geometry regularization.

\subsection{B. Detached transition-pooled conditional surrogate}
\begin{equation}
g_{\mathrm{cond}}
=\mu_0\nabla B_0+\mu_1\nabla B_1,
\label{eq:conditional}
\end{equation}
where the $\mu_s$ values are held fixed during differentiation.  This equals
the population expectation of computing a categorical barrier gradient at a
state drawn from the transition-pooled dataset.  It directly protects each
state's categorical action geometry in proportion to that state's dataset
frequency.  Consequently, direct protection of rarely visited states is weaker.

\subsection{C. Complete transition-pooled-weighted scalar objective}
\begin{equation}
B_{\mathrm{weighted}}=\mu_0B_0+\mu_1B_1.
\end{equation}
Its complete derivative is
\begin{equation}
\boxed{
\nabla B_{\mathrm{weighted}}
=g_{\mathrm{cond}}+B_0\nabla\mu_0+B_1\nabla\mu_1.}
\label{eq:weighted-gradient}
\end{equation}
This method asks what changes when we differentiate the policy-dependent state
weights instead of detaching them.  The extra term is not a categorical action
barrier; it comes from changing which states populate the pooled dataset.

\subsection{D. Uniform-state action barrier}
\begin{equation}
B_{\mathrm{uniform}}=\frac12(B_0+B_1).
\end{equation}
This is an oracle fixed-state-reference baseline.  It protects both statewise
categorical policies equally, regardless of how often the current policy visits
the states.  It is not directly available from ordinary on-policy sampled-state
averaging when a state is absent.

\subsection{E. Visitation-only barrier}
From Equation~\eqref{eq:mu},
\begin{equation}
B_{\mathrm{visit}}
=\log\mu_0+\log\mu_1
=\log q-2\log(1+q).
\end{equation}
Therefore
\begin{equation}
\frac{dB_{\mathrm{visit}}}{dq}
=\frac{1-q}{q(1+q)}>0,
\qquad 0<q<1.
\label{eq:visit-derivative}
\end{equation}
It always applies explicit pressure to visit $s_1$, independently of reward and
independently of the quality of $\pi_1$.

\subsection{F. Normalized full pooled-Fisher barrier}
\begin{equation}
B_{\mathrm{full}}
=B_{\mathrm{uniform}}+B_{\mathrm{visit}},
\qquad
\nabla B_{\mathrm{full}}
=\nabla B_{\mathrm{uniform}}+\nabla B_{\mathrm{visit}}.
\label{eq:full-gradient}
\end{equation}
This is the exact normalized log-determinant barrier for the declared
transition-pooled Fisher.  It combines equal statewise action protection with
explicit visitation pressure.

\subsection{The three objects that must not be confused}

Because $\mu_\theta$ depends on the policy,
\begin{equation}
\underbrace{\E_{s\sim\mu_\theta}[\nabla B_s]}_{\text{detached conditional}}
\neq
\underbrace{\nabla\E_{s\sim\mu_\theta}[B_s]}_{\text{complete weighted}}
\neq
\underbrace{\nabla\log\det F_{\pool}}_{\text{full pooled geometry}}
\end{equation}
in general.  Step 3 exists primarily to make this distinction explicit.

\section{Exact implementation and verification philosophy}

\subsection{Why there is no sampling in this stage}

Sampling would mix three questions: whether the algebra is correct, whether a
finite-batch estimator is biased, and whether its variance affects training.
This stage isolates the first question and the population dynamics.  It uses:
\begin{itemize}
\item CPU float64;
\item exact action and trajectory enumeration;
\item analytical reward and barrier gradients;
\item no Gym, neural network, Monte Carlo sampling, damping, clipping,
spectral floor, or stochastic noise.
\end{itemize}
Therefore each curve is one deterministic trajectory.  Seeds and confidence
intervals are neither needed nor meaningful here.

\subsection{Independent checks}

The verification CLI checks:
\begin{itemize}
\item trajectory probabilities and exact return enumeration;
\item the population transition counts underlying $\mu_0,\mu_1$;
\item zero expected scores and score-enumerated statewise Fishers;
\item the enumerated pooled Fisher and block formula;
\item symmetry, positive definiteness, and positive \texttt{slogdet} signs;
\item Equation~\eqref{eq:pooled-logdet};
\item the analytical reward gradient against exact REINFORCE enumeration;
\item detached, complete-weighted, visitation, and full gradients against
autograd and central finite differences;
\item Equations~\eqref{eq:weighted-gradient} and \eqref{eq:full-gradient};
\item deterministic replay, probability constructors, resume behavior, and
artifact schemas.
\end{itemize}
All checks pass at the declared float64 tolerances.

\subsection{Code objects and responsibilities}

\begin{center}\small
\begin{tabularx}{\textwidth}{lX}
\toprule
Object/module & Responsibility\\ \midrule
\texttt{TwoStepTrap} & MDP constants, trajectory enumeration, exact return,
and exact reward gradient.\\
\texttt{model.py} & Reduced logits, policies, $q$, $\mu_0,\mu_1$, and stable
probability-to-logit constructors.\\
\texttt{geometry.py} & Scores, statewise and pooled Fishers, all barrier values,
analytical gradients, and geometry snapshots.\\
\texttt{ExactTrainingConfig} & Method, $\alpha$, $\beta$, update count, and
experiment label.\\
\texttt{ExactTrainingResult} & Complete parameter and metric histories plus
finite-status flags; compressed deterministic serialization.\\
\texttt{reporting.py} & CSV summaries, trajectory figures, gradient
decompositions, phase plots, and basin maps.\\
\texttt{verify.py} & Independent algebraic, autograd, finite-difference, and
short robustness checks.\\
\texttt{run\_experiment.py} & Main, magnitude-controlled, robustness, and basin
suites; manifest validation and resume.\\ \bottomrule
\end{tabularx}
\end{center}

\section{Experiment 1: main same-beta comparison}

\subsection{Design and meaning}

The declared setting is
\begin{equation}
\alpha=0.05,qquad N_{\mathrm{updates}}=2000,
\qquad\beta\in\{0.01,0.1,0.2\}.
\end{equation}
All methods are first compared at the same numerical $\beta$.  This is the
primary comparison because they are genuinely different regularized
objectives, not approximations forced to have the same gradient norm.

The two initializations are:
\begin{align}
\phi_{\mathrm{uniform}}&=(0,0,0,0),\\
\phi_{\mathrm{adverse}}&=(2,-2,-2,2).
\end{align}
Uniform gives $q=\pi_1(a_0)=1/3$.  Adverse gives approximately
$q=\pi_1(a_0)=0.0159$: $s_0$ strongly prefers the safe action, while $s_1$
strongly prefers the mediocre action.  This is a predeclared difficult basin,
not a seed selected after observing results.

$\beta=0.01$ tests a weak barrier, $\beta=0.1$ tests an intermediate scale, and
$\beta=0.2$ tests whether stronger regularization creates visible bias.  This
parallels the barrier-strength trade-off emphasized by Cesani et al.
\cite[Secs.~6.1 and E.2/E.5]{cesani2026}, although our deterministic MDP is not
their stochastic Gaussian bandit experiment.

\subsection{Main adverse endpoint}

\begin{center}\small
\begin{tabular}{lrrr}
\toprule Method, $\beta=0.1$ & $J_T$ & $q_T$ & $\pi_1(a_0)_T$\\ \midrule
Reward only & .491 & .002 & .016\\
Detached conditional & .401 & .248 & .309\\
Complete weighted & .387 & .231 & .233\\
Uniform action & .565 & .533 & .699\\
Visitation only & .419 & .192 & .017\\
Full pooled Fisher & \textbf{.785} & \textbf{.830} & \textbf{.866}\\ \bottomrule
\end{tabular}
\par\smallskip\textit{Table 1: Exact endpoints from the adverse initialization.}
\end{center}

Reward-only ascent remains in the safe-action basin.  Visitation-only raises
$q$ but does not repair $\pi_1$.  Uniform action repairs $\pi_1$ without the
same direct pressure to reach it.  Full combines both and obtains the largest
return in this declared comparison.

From the uniform initialization, reward-only reaches $J=0.973$ and full at
$\beta=0.1$ reaches $0.795$.  Thus regularization helps the adverse recovery
but biases a policy that already follows a favorable path.  At the adverse
start, $\beta=0.01$ is weak.  Increasing full to $\beta=0.2$ reduces its return
from $0.785$ to $0.621$, while some conditional methods need the larger force
to escape.  ``Stronger barrier'' is therefore not a universal performance order.

\begin{figure}[htbp]\centering
\includegraphics[width=.94\textwidth]{plots/main/adverse_beta_0p1/performance_policy.png}
\caption{Main adverse trajectories.  Return alone is insufficient: $q$ says
whether $s_1$ is reached, while $\pi_1(a_0)$ says whether the policy behaves
well once it arrives.}
\label{fig:performance}
\end{figure}

\begin{figure}[htbp]\centering
\includegraphics[width=.86\textwidth]{plots/main/adverse_beta_0p1/phase_q_vs_good.png}
\caption{Adverse policy-space phase plot.  Horizontal movement primarily
changes visitation; vertical movement primarily repairs the good action at
$s_1$.  Full moves along both coordinates.}
\label{fig:phase}
\end{figure}

\section{Experiment 2: magnitude-controlled comparison}

\subsection{What problem this control addresses}

Equal $\beta$ does not imply equal update size because the five barrier
gradients have different norms.  A method could appear stronger only because
$\|\beta\nabla B\|$ is initially larger.  At the adverse initialization we set
\begin{equation}
\tau_{\mathrm{mag}}=\|\nabla J(\phi_0)\|_2,
\qquad
\beta_m=\frac{\tau_{\mathrm{mag}}}{\|\nabla B_m(\phi_0)\|_2}.
\label{eq:magnitude-match}
\end{equation}
This makes the initial reward and regularizer contributions equal in norm.
Reward-only remains an unscaled reference.  This is a secondary diagnostic,
not a replacement for the main same-$\beta$ objectives.

The resulting coefficients are $0.0301$ (detached and complete), $0.0420$
(uniform), $0.0435$ (visitation), and $0.0234$ (full).

\begin{center}\small
\begin{tabular}{lrrr}
\toprule Method & $J_T$ & $q_T$ & $\pi_1(a_0)_T$\\ \midrule
Reward only & .491 & .002 & .016\\
Detached conditional & .443 & .057 & .021\\
Complete weighted & .444 & .055 & .021\\
Uniform action & .459 & .039 & .242\\
Visitation only & .453 & .078 & .016\\
Full pooled Fisher & .453 & .069 & .135\\ \bottomrule
\end{tabular}
\par\smallskip\textit{Table 2: Endpoints after initial magnitude matching.}
\end{center}

No regularizer recovers within 2000 updates under this matched target.
Consequently, full's strong same-$\beta$ recovery partly depends on scale.
Direction still matters: uniform mainly changes $\pi_1$, visitation mainly
changes $q$, and full changes both.

\begin{figure}[htbp]\centering
\includegraphics[width=.94\textwidth]{plots/main/adverse_beta_0p1/gradient_decompositions.png}
\caption{Component norms and cosines.  The top row decomposes full into
uniform-action and visitation directions.  The bottom row decomposes complete
weighted into the detached conditional direction and the derivative of the
state weights.}
\label{fig:decomposition}
\end{figure}

\section{Experiment 3: explicit-Euler step-size robustness}

The main result uses $\alpha=0.05$.  To test whether its qualitative behavior
is a discretization artifact, the robustness study uses
\begin{equation}
\alpha\in\{0.025,0.05,0.1\}
\end{equation}
while keeping optimization time $\alpha N=100$ fixed.  The corresponding
update counts are $4000$, $2000$, and $1000$.  Holding the update count fixed
would change both the step size and total optimization time, so it would not
isolate Euler discretization.

All 36 runs remain finite.  Within a fixed method and initialization, the
largest endpoint-return range across step sizes is $1.1\times10^{-4}$.
Endpoint ordering and phase behavior are unchanged.  This supports the narrow
claim that the reported exact dynamics are not artifacts of choosing only
$\alpha=0.05$; it is not a hyperparameter search or a stability proof for
arbitrary step sizes.

\section{Experiment 4: deterministic initialization-basin grid}

\subsection{Constructor and purpose}

Two scalar probabilities are varied:
\begin{equation}
q_0\in\mathcal G,qquad \good{}_0=\pi_1(a_0)\in\mathcal G,
\quad
\mathcal G=\{.02,.05,.1,.25,.5,.75,.9\}.
\end{equation}
The remaining probability is split in a fixed, symmetric way:
\begin{align}
\pi_0&=(\tfrac12(1-q_0),q_0,\tfrac12(1-q_0)),\\
\pi_1&=(\good{}_0,\tfrac12(1-\good{}_0),\tfrac12(1-\good{}_0)).
\end{align}
For reference action $a_2$, the stable reduced-logit constructor uses
\begin{equation}
\phi_a=\log\frac{\pi(a)}{\pi(a_2)}.
\end{equation}
The grid asks whether conclusions based on only uniform and adverse starts are
localized accidents or reflect broader basins.

Each map reports final return, final $q$, final good-action probability, and a
visualization-only classification
\begin{equation}
q_T\geq0.9\quad\text{and}\quad\pi_1(a_0)_T\geq0.9.
\end{equation}
This threshold is not a test assertion and does not define the scientific
winner.  A regularized interior equilibrium can have good return without
crossing $0.9$.

At $\beta=0.1$, full has the highest worst-case return ($0.748$) and mean return
($0.789$), but zero points cross the displayed threshold because its equilibrium
remains interior.  Reward only crosses at 20/49 points; visitation only at
33/49.  These counts describe phase behavior, not general superiority.

\begin{figure}[p]\centering
\begin{minipage}{.48\textwidth}\includegraphics[width=\textwidth]{plots/basin/reward_only.png}\end{minipage}
\begin{minipage}{.48\textwidth}\includegraphics[width=\textwidth]{plots/basin/detached_conditional.png}\end{minipage}
\caption{Basin maps for reward-only and detached conditional training.}
\end{figure}
\begin{figure}[p]\centering
\begin{minipage}{.48\textwidth}\includegraphics[width=\textwidth]{plots/basin/complete_weighted.png}\end{minipage}
\begin{minipage}{.48\textwidth}\includegraphics[width=\textwidth]{plots/basin/uniform_action.png}\end{minipage}
\caption{Basin maps for complete-weighted and uniform-action training.}
\end{figure}
\begin{figure}[p]\centering
\begin{minipage}{.48\textwidth}\includegraphics[width=\textwidth]{plots/basin/visitation_only.png}\end{minipage}
\begin{minipage}{.48\textwidth}\includegraphics[width=\textwidth]{plots/basin/full_pooled_fisher.png}\end{minipage}
\caption{Basin maps for visitation-only and full pooled-Fisher training.}
\end{figure}

\section{Diagnostics: what every recorded quantity means}

At every update the archive records the complete $\pi_0$, $\pi_1$, $q$,
$\mu_0$, $\mu_1$, and the following groups.

\subsection{Performance and behavioral coordinates}
\begin{itemize}
\item $J$: exact expected return.
\item $q=\pi_0(a_1)$: probability of reaching $s_1$.
\item $\pi_1(a_0)$: probability of choosing the good action after reaching it.
\item $\min_a\pi_s(a)$: direct statewise action-support diagnostic.
\end{itemize}

\subsection{Statewise and pooled geometry}
\begin{itemize}
\item $\lambda_{\min}(F_0)$ and $\lambda_{\min}(F_1)$: weakest conditional
categorical directions.
\item $\lambda_{\min}(F_{\pool})$: weakest direction after both conditional
geometry and dataset frequency are combined.
\item $\log\det F_0$, $\log\det F_1$, $\log\det F_{\pool}$: statewise and pooled
log-volume measures.
\item $B_0$, $B_1$, $B_{\mathrm{weighted}}$, $B_{\mathrm{uniform}}$,
$B_{\mathrm{visit}}$, $B_{\mathrm{full}}$: the scalar objectives needed to
audit every decomposition.
\end{itemize}

\begin{figure}[htbp]\centering
\includegraphics[width=.82\textwidth]{plots/main/adverse_beta_0p1/statewise_geometry.png}
\caption{Statewise minima and Fisher eigenvalues.  These identify whether a
pooled change is caused by conditional degeneration or by reduced visitation.}
\label{fig:statewise}
\end{figure}

\begin{figure}[htbp]\centering
\includegraphics[width=.90\textwidth]{plots/main/adverse_beta_0p1/pooled_statewise_logdet.png}
\caption{Pooled and statewise log determinants.  The pooled curve alone cannot
identify which state or weighting mechanism changed.}
\label{fig:logdet}
\end{figure}

\begin{figure}[htbp]\centering
\includegraphics[width=.88\textwidth]{plots/main/adverse_beta_0p1/action_visitation_barriers.png}
\caption{The normalized full barrier separates exactly into an action term and
a visitation term.  Plotting both is required to interpret the full update.}
\label{fig:barriers}
\end{figure}

\subsection{Gradient scale and direction}

Every method records
\begin{equation}
\|\nabla J\|_2,qquad
\|\beta\nabla B\|_2,qquad
\cos(\nabla J,\nabla B).
\end{equation}
The norm comparison says which component dominates update magnitude.  The
cosine says whether the regularizer locally supports the reward direction
($+1$), opposes it ($-1$), or is approximately orthogonal ($0$).  A zero-norm
case is stored as undefined/NaN, not stabilized with an arbitrary epsilon.

For full, the uniform and visitation components are recorded separately.  For
complete weighted, the detached and state-weight-derivative components are
also separated.  These diagnostics make the magnitude-controlled conclusion
possible.

\section{What the exact results support}

\subsection{Supported statements}

\begin{itemize}
\item The detached transition-pooled conditional surrogate directly protects
each state's categorical action geometry in proportion to that state's
frequency in the rollout dataset.
\item Consequently, its direct protection of rarely visited states is
proportionally weaker.
\item The complete derivative of the transition-pooled-weighted scalar
objective additionally contains derivatives of the policy-dependent state
weights.
\item The full pooled-Fisher log determinant contains explicit visitation
pressure in addition to the fixed-state conditional action barrier.
\item In this MDP and adverse initialization, combining these directions
improves recovery at the declared same $\beta$.
\item The same barrier biases the solution away from the reward-optimal
deterministic policy, agreeing with the exploration/bias trade-off emphasized
by Cesani et al. \cite[Sec.~6.1 and App.~E.2]{cesani2026}.
\item The magnitude-controlled experiment shows that scale explains part of
the same-$\beta$ difference; the component trajectories show that direction
also differs.
\end{itemize}

\subsection{Claims this experiment does not support}

It does not establish that:
\begin{itemize}
\item the detached conditional surrogate equals a global neural Fisher log
determinant;
\item a larger determinant guarantees better conditioning or return;
\item improved return proves Fisher geometry caused the improvement;
\item the six objectives should share an optimal numerical $\beta$;
\item this two-state result transfers directly to shared neural parameters;
\item the empirical finite-batch Fisher is interchangeable with the population
Fisher or the Hessian \cite{Kunstner2019}.
\end{itemize}

The exact result is a mechanism isolation: it shows which derivative terms
exist and how they move this controlled MDP.

\section{Artifacts and exact reconstruction procedure}

From the repository root, first verify the mathematics:
\begin{center}\footnotesize\ttfamily
python -m exploration.tabular\_mdp.verify
\end{center}
Then run or resume the complete deterministic suite:
\begin{center}\footnotesize\ttfamily
python -m exploration.tabular\_mdp.run\_experiment\\
python -m exploration.tabular\_mdp.run\_experiment --resume
\end{center}

The output root is
\path{exploration/results/tabular_mdp/two_step_trap/}.  Its main objects are:
\begin{itemize}
\item \texttt{config.json}: complete, versioned experiment declaration;
\item \texttt{verification.json}: residuals and pass/fail decisions;
\item \texttt{magnitude\_matched\_betas.json}: computed secondary coefficients;
\item \texttt{archives/}: 80 compressed exact trajectory units;
\item \texttt{summary.csv} and \texttt{run\_status.csv}: endpoint and finite
status tables;
\item \texttt{plots/}: 97 figures covering main, magnitude, robustness, and
basin analyses.
\end{itemize}
The output manifest is checked before resume.  A nonempty incompatible output
directory is rejected rather than overwritten.

\section{Decision gate and next experiment}

The next step should retain the same two-state MDP but sample trajectories.
It should compare population quantities against finite batches and measure:
\begin{itemize}
\item bias and variance of the sampled gradients and pooled Fisher;
\item the random-denominator effect in Equation~\eqref{eq:ratio-estimator};
\item the frequency of batches containing zero $s_1$ transitions;
\item how those empty batches weaken the sampled-state conditional barrier;
\item whether increasing batch size approaches the exact Step 3 dynamics.
\end{itemize}

Only afterward should the work move to shared-parameter categorical neural
policies such as CartPole and Acrobot.  The practical method should be called
an \emph{on-policy sampled-state conditional categorical log barrier}, not the
global neural Fisher log determinant.  Continuous Gaussian policies remain a
separate extension.

\section{Source-reading guide}

For verification beyond this document:
\begin{center}\small
\begin{tabularx}{\textwidth}{lX}
\toprule
Question & Recommended source location\\ \midrule
Why start with MAB? & Cesani et al., Sections 1--3.\\
Why is the categorical barrier a Fisher log determinant? & Cesani et al.,
Section 4 and Appendix B.3.2.\\
What update does LB-SGB actually use? & Cesani et al., Section 5 and Appendix
B.4, especially the gradient $\mathbf1-K\pi$.\\
Where is the exploration/bias trade-off? & Cesani et al., Section 6.1 and
Appendix E.2/E.5.\\
Where is the warning about general RL? & Cesani et al., Section 8.\\
Why does the Fisher define a natural policy metric? & Kakade, Sections 2--3;
Amari's original natural-gradient article.\\
Why is Equation~\eqref{eq:reinforce} a policy gradient? & Williams's REINFORCE
paper and Sutton et al.'s policy-gradient theorem.\\
Why distinguish empirical and population Fisher? & Kunstner et al., especially
Sections 2--4.\\ \bottomrule
\end{tabularx}
\end{center}

\begin{sloppypar}
\begin{thebibliography}{9}

\bibitem{cesani2026}
L. Cesani, M. Papini, and M. Restelli,
``How Log-Barrier Helps Exploration in Policy Optimization,''
arXiv:2603.15001v2, 2026.
\url{https://arxiv.org/html/2603.15001v2}

\bibitem{Williams1992}
R. J. Williams,
``Simple Statistical Gradient-Following Algorithms for Connectionist
Reinforcement Learning,''
\emph{Machine Learning}, 8, pp. 229--256, 1992.
\url{https://doi.org/10.1007/BF00992696}

\bibitem{SuttonMcAllester2000}
R. S. Sutton, D. McAllester, S. Singh, and Y. Mansour,
``Policy Gradient Methods for Reinforcement Learning with Function
Approximation,'' \emph{NeurIPS 12}, 2000.
\href{https://papers.nips.cc/paper/1999/hash/464d828b85b0bed98e80ade0a5c43b0f-Abstract.html}{Official NeurIPS proceedings page.}

\bibitem{Kakade2001}
S. Kakade,
``A Natural Policy Gradient,'' \emph{NeurIPS 14}, 2001.
\href{https://papers.nips.cc/paper/2001/hash/4b86abe48d358ecf194c56c69108433e-Abstract.html}{Official NeurIPS proceedings page.}

\bibitem{Amari1998}
S.-i. Amari,
``Natural Gradient Works Efficiently in Learning,''
\emph{Neural Computation}, 10(2), pp. 251--276, 1998.
\url{https://doi.org/10.1162/089976698300017746}

\bibitem{Kunstner2019}
F. Kunstner, P. Hennig, and L. Balles,
``Limitations of the Empirical Fisher Approximation for Natural Gradient
Descent,'' \emph{NeurIPS 32}, 2019.
\href{https://proceedings.neurips.cc/paper/2019/hash/46a558d97954d0692411c861cf78ef79-Abstract.html}{Official NeurIPS proceedings page.}

\bibitem{agarwal2020}
A. Agarwal, M. Henaff, S. Kakade, and W. Sun,
``PC-PG: Policy Cover Directed Exploration for Provable Policy Gradient
Learning,'' \emph{NeurIPS 33}, 2020.
\href{https://papers.nips.cc/paper/2020/hash/9b3a9fb4db30fc6594ec3990cbc09932-Abstract.html}{Official NeurIPS proceedings page.}

\end{thebibliography}
\end{sloppypar}

\end{document}
